\chapter{طراحی و پیاده‌سازی}
\label{ch:design}

این فصل مقیاس‌گذار پیاده‌سازی‌شده را شرح می‌دهد. ترتیب ارائه از اصول طراحی آغاز می‌شود، سپس
هر مؤلفه جداگانه توضیح داده می‌شود، و در پایان نقص‌هایی که در جریان پیاده‌سازی آشکار شدند ثبت
می‌گردد.

دو تصمیم طراحی در این فصل، انحراف عمدی از نمونه‌ی اولیه‌ی پروژه‌اند و هر دو با اندازه‌گیری
توجیه شده‌اند: پایه‌ی فرمول باید نمونه‌های \emph{آماده} باشد (بخش~\ref{sec:design-formula})،
و پیش‌بینی باید روی بار \emph{کل} انجام شود (بخش~\ref{sec:design-perreplica}). تصمیم دوم،
هسته‌ی علمی این پروژه است.

\section{اصول طراحی}
\label{sec:design-principles}

پیشنهاد پروژه اصل «تفکیک دغدغه‌ها»\LTRfootnote{Separation of Concerns} را مبنا گرفته بود. بر
همین پایه، سامانه در امتداد سه پرسش مستقل شکسته شده است: یک سیگنال بخوان، یک اندازه‌ی ناوگان
تصمیم بگیر، آن را اعمال کن. هر یک از این سه مرحله یک واسط است و هر واسط دست‌کم دو پیاده‌سازی
دارد.

این تصمیم سلیقه‌ای نیست، چون یک پیامد عملی مشخص دارد. منطق کنترل بدون \lr{Prometheus} و بدون
خوشه‌ی \lr{Kubernetes} آزمودنی می‌شود. کافی است جمع‌آورنده‌ی ساختگی و مقیاس‌گذار
درون‌حافظه‌ای را جای پیاده‌سازی‌های واقعی بگذاریم تا یک حلقه‌ی کنترل \emph{بسته} به‌طور کامل
درون یک آزمون واحد اجرا شود. آنگاه می‌توان دینامیک سیاست را مستقل از هر زیرساختی سنجید.
چنان‌که در بخش~\ref{sec:design-perreplica} خواهیم دید، مهم‌ترین یافته‌ی این پروژه از همین
امکان به دست آمد.

سه اصل دیگر نیز در سراسر پیاده‌سازی رعایت شده است:

\begin{itemize}
  \item \textbf{هیچ رفتاری در کد سخت‌کد نشود که بتواند پیکربندی باشد.} عبارت \lr{PromQL}،
  انتخاب سیاست، ضرایب پیش‌بینی‌کننده و پارامترهای پایدارسازی همگی از یک فایل \lr{YAML}
  می‌آیند. این همان نیازمندی کارکردی «پشتیبانی از سیاست‌های چندگانه و پیکربندی‌پذیر بدون
  تغییر در کد» است.
  \item \textbf{در نبود اطلاعات، هیچ کاری نکن.} شکست در هر مرحله از حلقه، تعداد نمونه‌ی جاری
  را دست‌نخورده می‌گذارد. دلیل آن در بخش~\ref{sec:design-loop} آمده است.
  \item \textbf{تصمیم باید خودش را توضیح دهد.} ساختار \lr{Decision} علاوه بر پاسخ، ورودی‌ها و
  دلیل را نیز حمل می‌کند، و همه‌ی این‌ها هم در گزارش‌های ساختاریافته و هم روی نقطه‌ی
  \lr{/metrics} منتشر می‌شود.
\end{itemize}

پیاده‌سازی به زبان \lr{Go} است و حدود ۱۹۸۵ خط کد غیرآزمونی به‌همراه حدود ۲۰۲۰ خط کد آزمون
دارد.

\section{معماری سامانه}
\label{sec:design-architecture}

شکل~\ref{fig:architecture} معماری سامانه را نشان می‌دهد و همان نمودار بلوکی صفحه‌ی ۵ پیشنهاد
پروژه را بازتولید می‌کند.

\begin{figure}[H]
\centering
\begin{latin}
\resizebox{\linewidth}{!}{%
\begin{tikzpicture}[
  node distance=8mm,
  box/.style={draw, rounded corners=2pt, align=center, minimum height=9mm,
              inner sep=4pt, font=\small},
  ext/.style={box, fill=black!5},
  core/.style={box, fill=blue!8},
  app/.style={box, fill=orange!12},
  >={Stealth[length=2mm]},
  every path/.style={->, thick}
]
  \node[ext]  (lg)   {Load Generator\\(k6)};
  \node[app]  (app)  [right=16mm of lg] {Target Microservice\\(sample app)};
  \node[ext]  (prom) [right=16mm of app] {Prometheus\\(metrics + PromQL)};
  \node[ext]  (graf) [right=14mm of prom] {Grafana\\Dashboard};

  \node[core] (mc)   [below=18mm of app]  {Metric Collector};
  \node[core] (pe)   [right=10mm of mc]   {Policy Engine\\(threshold / predictive)};
  \node[core] (sc)   [right=10mm of pe]   {Scaler\\(Reconciler)};

  \node[app]  (api)  [below=14mm of mc]  {Kubernetes API\\(scale subresource)};

  \draw (lg)   -- node[above, font=\tiny] {traffic}  (app);
  \draw (app)  -- node[above, font=\tiny] {scrape}   (prom);
  \draw (prom) -- (graf);
  \draw (prom.south) -- node[right, font=\tiny] {PromQL} (mc.north);
  \draw (mc)   -- node[above, font=\tiny] {metric}   (pe);
  \draw (pe)   -- node[above, font=\tiny] {decision} (sc);
  \draw (sc.south) |- node[near end, above, font=\tiny] {patch replicas} (api.east);
  \draw (api.west) -- ++(-10mm,0) |- (app.west);

  \begin{scope}[on background layer]
    \node[draw, dashed, rounded corners=3pt, fit=(mc)(pe)(sc),
          inner sep=6mm, label={[font=\small\bfseries]below:Custom Autoscaler (Go)}] {};
  \end{scope}
\end{tikzpicture}%
}
\end{latin}
\caption[معماری سامانه]{معماری سامانه. مقیاس‌گذار سفارشی (کادر نقطه‌چین) سیگنال را از
\lr{Prometheus} می‌خواند، اندازه‌ی ناوگان را تصمیم می‌گیرد، و آن را از راه زیرمنبع
\lr{scale} اعمال می‌کند. سنجه‌های داخلی خود مقیاس‌گذار نیز به \lr{Prometheus} داده می‌شود.}
\label{fig:architecture}
\end{figure}

مؤلفه‌ها مطابق جدول~\ref{tab:design-components} به بسته‌های کد نگاشته می‌شوند:

\begin{table}[H]
\centering
\small
\caption[نگاشت مؤلفه‌های معماری به بسته‌های کد]{نگاشت مؤلفه‌های معماری به بسته‌ها و
 واسط‌های کد. همه‌ی بسته‌ها زیر \code{internal/} قرار دارند.}
\label{tab:design-components}
\scriptsize
\begin{tabular}{|l|l|l|l|}
\hline
\textbf{مؤلفه} & \textbf{بسته} & \textbf{واسط} & \textbf{پیاده‌سازی‌ها} \\
\hline
جمع‌آورنده‌ی سنجه & \code{metrics} & \code{Collector.Read} & \code{Prometheus}، \code{Static} \\
موتور سیاست & \code{policy} & \code{Policy.Decide} & \code{Threshold}، \code{PredictiveTotalLoad}، \\
 & & & \code{PredictivePerReplica} \\
پیش‌بینی‌کننده & \code{policy} & \code{Forecaster.Update} & \code{EWMATrend}، \code{Holt} \\
مقیاس‌گذار & \code{scaler} & \code{Scaler.Get} / \code{Set} & \code{Kubernetes}، \code{InMemory} \\
آشتی‌دهنده & \code{controller} & \lr{---} & حلقه‌ی کنترل \\
صادرکننده‌ی سنجه & \code{observability} & \lr{---} & نقطه‌ی \code{/metrics} \\
\hline
\end{tabular}
\end{table}

ورودی سامانه سنجه‌های سطح‌برنامه از \lr{Prometheus} است: نرخ درخواست، تأخیر، طول صف. این
سنجه‌ها از راه یک عبارت \lr{PromQL} خوانده می‌شوند که در پیکربندی تعیین می‌شود. خروجی سامانه
تعداد نمونه‌ی \lr{Deployment} هدف است، به‌علاوه‌ی سنجه‌های داخلی خود مقیاس‌گذار برای
مصورسازی در \lr{Grafana}.

\section{جمع‌آورنده‌ی سنجه}
\label{sec:design-collector}

جمع‌آورنده یک عبارت \lr{PromQL} را از پیکربندی می‌خواند و اجرا می‌کند. این عبارت در کد
سخت‌کد نشده است. همین نکته ادعای «مقیاس‌گذاری بر پایه‌ی سنجه‌های سطح‌برنامه» را از یک ادعای
کلی به یک قابلیت واقعی تبدیل می‌کند. نرخ درخواست، تأخیر یا طول صف تنها با تغییر یک رشته در
فایل پیکربندی جای یکدیگر را می‌گیرند و کد دست‌نخورده می‌ماند.

عبارتی که در ارزیابی به کار رفت:

\begin{latin}
\begin{lstlisting}[language=promql]
sum(rate(http_requests_total{app="sample"}[1m]))
  / clamp_min(count(up{app="sample"} == 1), 1)
\end{lstlisting}
\end{latin}

در همین عبارت کوتاه دو تصمیم طراحی نهفته است. نخست آنکه مخرج، تعداد نمونه‌هایی است که
\lr{Prometheus} آن‌ها را در دسترس می‌بیند و نه \lr{spec.replicas}. این انتخاب سیگنال را با
پایه‌ی «نمونه‌ی آماده» هم‌خوان نگه می‌دارد که سیاست به کار می‌برد و در
بخش~\ref{sec:design-formula} توضیح داده می‌شود. دوم آنکه \lr{clamp\_min} جلوی تقسیم بر صفر را
می‌گیرد، برای لحظه‌ای که هیچ نمونه‌ای آماده نیست.

نکته‌ی سومی نیز هست که در فصل~\ref{ch:results} اهمیت پیدا می‌کند: صورت کسر یک \lr{sum} روی
\emph{همه‌ی} سری‌هاست و مخرج تعداد هدف‌های در دسترس را می‌شمارد. پس هیچ‌یک از دو طرف به این
وابسته نیست که \emph{کدام} نمونه‌ها ترافیک گرفته‌اند. همین ویژگی است که باعث می‌شود نقص تحویل
بارِ توضیح‌داده‌شده در بخش~\ref{sec:results-loaddefect}، ورودی مقیاس‌گذار را خراب نکند.

رفتار در خطا هم صریح است. اگر خواندن سنجه ناموفق بماند، جمع‌آورنده خطا برمی‌گرداند و حلقه‌ی
کنترل تعداد نمونه‌ی فعلی را نگه می‌دارد. منطق این تصمیم ساده است: نبود سنجه شاهدی بر تغییر بار
در هیچ‌یک از دو جهت نیست، پس امن‌ترین کار هیچ‌کاری‌نکردن است.

پیاده‌سازی دوم این واسط، \code{Static}، مقداری ثابت یا دنباله‌ای از پیش تعیین‌شده برمی‌گرداند
و همان چیزی است که آزمون‌های حلقه‌بسته را ممکن می‌کند.

\section{موتور سیاست}
\label{sec:design-policy}

\subsection{فرمول پایه: نمونه‌های آماده، نه نمونه‌های درخواست‌شده}
\label{sec:design-formula}

تعداد نمونه‌ی هدف از همان فرمولی می‌آید که \lr{HPA} به کار می‌برد
(برابری~\ref{eq:hpa} در فصل~\ref{ch:background}):

\begin{equation}
\label{eq:desired}
d = \left\lceil r \times \frac{m}{T} \right\rceil
\end{equation}

که در آن $r$ تعداد نمونه‌های آماده، $m$ سنجه‌ی مشاهده‌شده و $T$ نرخ هدف است.

نکته‌ی ظریف در انتخاب پایه است. نمونه‌ی اولیه‌ی پایتونی این پروژه \lr{spec.replicas} را در
فرمول می‌گذاشت، در حالی که سنجه به‌ازای هر نمونه‌ی \emph{آماده} اندازه گرفته می‌شد. این دو در
حالت پایدار یکی هستند اما تا وقتی نمونه‌های تازه در حال آماده‌شدن‌اند از هم واگرا می‌شوند.
همان بازه، دقیقاً پنجره‌ای است که در یک مقیاس‌گیری به بالا اهمیت دارد.

یک مثال عددی کافی است، که اکنون به شکل آزمون \code{TestFormulaBaseIsReadyNotSpec} تثبیت شده
است. فرض کنید $\mathrm{spec} = 4$، $r = 2$، $m = 200$ و $T = 100$. پاسخ درست ۴ است. اگر
\lr{spec} را پایه بگیریم پاسخ ۸ می‌شود، یعنی دو برابر. مقیاس‌گذار بار مشاهده‌شده را بر
نمونه‌هایی تقسیم می‌کند که درخواست کرده، نه نمونه‌هایی که واقعاً سرویس می‌دهند. پس
بیش‌ازاندازه مقیاس می‌گیرد.

به همین دلیل \code{Scaler} هر دو عدد را برمی‌گرداند و سیاست از \code{Ready} استفاده می‌کند.
این یک بهبود عمدی نسبت به نمونه‌ی اولیه است و نه صرفاً ترجمه‌ی آن.

نکته‌ی متقارن آن هم اهمیت دارد. تصمیم \emph{اعمال} با \lr{spec.replicas} مقایسه می‌شود و نه
با \lr{ready}. پرسش آن مرحله این است که آیا پیش‌تر همین تعداد از خوشه خواسته شده است یا نه.
نمونه‌هایی که هنوز بالا می‌آیند، خواسته شده‌اند. اگر با \lr{ready} مقایسه می‌کردیم، همان
مقیاس‌گیری به بالا در هر چرخه دوباره صادر می‌شد تا وقتی نمونه‌ها بالا بیایند.

\subsection{ناحیه‌ی مرده}
\label{sec:design-deadband}

مانند \lr{HPA}، تغییری اعمال نمی‌شود مگر آنکه نسبت $m/T$ بیش از یک آستانه‌ی نسبی $\varepsilon$
از یک فاصله بگیرد:

\begin{equation}
\label{eq:deadband}
\left| \frac{m}{T} - 1 \right| \le \varepsilon \;\Longrightarrow\; d = r
\end{equation}

این آستانه به‌صورت پیش‌فرض ۱۰ درصد است. ناحیه‌ی مرده جلوی تعقیب نوسان‌های کوچک سنجه را
می‌گیرد.

مرز دقیق این ناحیه عمداً در هیچ آزمونی تثبیت نشده است. دلیلش رفتار ممیز شناور است. مثلاً
\lr{110/100} برابر \lr{1.1000000000000001} ارزیابی می‌شود، پس $|m/T - 1|$ اندکی از ۱۰ درصد
بزرگ‌تر می‌شود و ناوگان حرکت می‌کند. جفت دیگری از سنجه و هدف با همان نسبت اسمی ممکن است
آن‌سوی مرز بیفتد و ثابت بماند. \lr{HPA} استاندارد نیز همین نسبت را به همین شکل حساب می‌کند و
همین ابهام را به ارث می‌برد. اگر این مرز را در آزمون تثبیت می‌کردیم، به‌جای یک تصمیم طراحی،
مصنوعی از نمایش دودویی اعداد را تثبیت کرده بودیم.

\subsection{سیاست آستانه‌ای}
\label{sec:design-threshold}

ساده‌ترین سیاست است. سنجه‌ی جاری را در فرمول~\ref{eq:desired} می‌گذارد و تعداد نمونه را
پیشنهاد می‌دهد. این سیاست کاملاً واکنشی است و در ارزیابی نقش خط مبنا را دارد. ارزش افزوده‌ی
پیش‌بینی در برابر همین خط سنجیده می‌شود.

\subsection{پیش‌بینی‌کننده‌ها}
\label{sec:design-forecasters}

دو پیش‌بینی‌کننده پشت واسط \code{Forecaster} پیاده‌سازی شده است. هر دو حالت‌دار و ارزان‌اند،
چون در هر چرخه‌ی کنترل یک بار فراخوانی می‌شوند.

\textbf{\code{EWMATrend}} \mdash{} پیش‌فرض پیکربندی. سری را با میانگین متحرک نمایی هموار می‌کند و
سپس اختلاف دو مقدار هموارشده‌ی متوالی را به‌عنوان برآورد شیب می‌گیرد و به‌صورت خطی برون‌یابی
می‌کند:

\begin{align}
\label{eq:ewmatrend}
\ell_k &= \alpha\, x_k + (1-\alpha)\, \ell_{k-1} \\
\hat{x}_{k+h} &= \max\!\left(0,\; \ell_k + h\,(\ell_k - \ell_{k-1})\right)
\end{align}

\textbf{\code{Holt}} \mdash{} هموارسازی نمایی مضاعف مطابق برابری~\ref{eq:holt}، که سطح و شیب را با
دو ضریب جداگانه نگه می‌دارد و بنابراین یک روند پایدار را به اندازه‌ی \code{EWMATrend} میرا
نمی‌کند.

دو نکته‌ی پیاده‌سازی در هر دو مشترک است. نخست، کف \emph{نامنفی} روی خروجی: پیش‌بینی منفی برای
نرخ ورود درخواست بی‌معناست و \mdash{} چنان‌که در بخش~\ref{sec:design-defects} می‌آید \mdash{} اگر به
فرمول برسد به یک مقیاس‌گیری به پایینِ ناشی از نقص اندازه‌گیری می‌انجامد. دوم، مسیر جداگانه‌ای
برای ورودی \lr{NaN} یا بی‌نهایت، که در عمل یعنی یک \lr{scrape} ازدست‌رفته.

پیش‌بینی‌کننده‌ی دوم صرفاً برای آن هست که نشان دهد روش پیش‌بینی واقعاً افزون‌پذیر است و در
واسط قفل نشده. این همان ادعایی است که در نیازمندی‌های کارکردی آمده بود: پشتیبانی از سیاست‌های
چندگانه و پیکربندی‌پذیر، بدون تغییر در کد.

\subsection{سیاست پیش‌بینانه}
\label{sec:design-predictive}

سیاست پیش‌بینانه به‌جای مقدار جاری، مقدار پیش‌بینی‌شده برای چند گام بعد را در فرمول می‌گذارد.
فاصله‌ی آشتی‌دهی ۱۵ ثانیه و افق ۳ گام است، پس حدود ۴۵ ثانیه زمان پیش‌دستی به دست می‌آید. این
بازه تقریباً هم‌اندازه‌ی زمان آماده‌شدن یک نمونه‌ی جدید است و همین است که پیش‌بینی را سودمند
می‌کند.

\subsection{پیش‌بینی بار کل در برابر بار هر نمونه}
\label{sec:design-perreplica}

این بخش هسته‌ی علمی پروژه است. تصمیمی را ثبت می‌کند که از اندازه‌گیری به دست آمد و نه از
مطالعه‌ی کد.

نمونه‌ی اولیه بار \emph{هر نمونه} را پیش‌بینی می‌کرد. اما بار هر نمونه برابر است با بار کل
تقسیم بر نمونه‌های آماده،

\begin{equation}
\label{eq:perreplica}
m(t) = \frac{\lambda(t)}{r(t)}
\end{equation}

و $r(t)$ همان کمیتی است که مقیاس‌گذار تعیین می‌کند. پس پیش‌بینی‌کننده سیگنالی را برون‌یابی
می‌کند که کنش‌های خودش آن را جابه‌جا می‌کند. حلقه بهره‌ی مثبت پیدا می‌کند، مطابق آنچه
شکل~\ref{fig:feedback} نشان می‌دهد.

\begin{figure}[H]
\centering
\begin{latin}
\begin{tikzpicture}[
  node distance=6mm,
  n/.style={draw, rounded corners=2pt, align=center, font=\small,
            minimum height=8mm, inner sep=4pt},
  bad/.style={n, fill=red!8},
  good/.style={n, fill=green!8},
  >={Stealth[length=2mm]},
  every path/.style={->, thick}
]
  % --- unstable loop
  \node[bad] (a) {scale up};
  \node[bad] (b) [right=10mm of a] {per-replica\\load falls};
  \node[bad] (c) [right=10mm of b] {forecast\\trends down};
  \node[bad] (d) [right=10mm of c] {scale down};

  \draw (a) -- (b);
  \draw (b) -- (c);
  \draw (c) -- (d);
  \draw (d.south) -- ++(0,-6mm) -| node[near start, below, font=\tiny]
        {per-replica load rises $\rightarrow$ forecast trends up} (a.south);

  % --- stable loop
  \node[good] (e) [below=22mm of a] {arrival rate\\$\lambda(t)$};
  \node[good] (f) [right=10mm of e] {forecast\\$\hat{\lambda}$};
  \node[good] (g) [right=10mm of f] {fleet\\$\lceil \hat{\lambda}/T \rceil$};
  \node[good] (h) [right=10mm of g] {scale};

  \draw (e) -- (f);
  \draw (f) -- (g);
  \draw (g) -- (h);
  \draw[dashed, ->] (h.north) -- node[right, font=\tiny] {no path back} ++(0,6mm);
\end{tikzpicture}
\end{latin}
\caption[پس‌خور مثبت در پیش‌بینی بار هر نمونه]{بالا: پیش‌بینی بار هر نمونه یک مسیر بسته با
بهره‌ی مثبت می‌سازد. پایین: پیش‌بینی نرخ ورود، که برون‌زاست، چنین مسیری ندارد و اندازه‌ی
ناوگان به تعداد نمونه‌ی جاری وابسته نیست.}
\label{fig:feedback}
\end{figure}

\textbf{چرا سیاست واکنشی روی همین سیگنال پایدار می‌ماند.} پرسش طبیعی این است که سیاست
آستانه‌ای هم همین $m$ درون‌زا را می‌خواند، پس چرا نوسان نمی‌کند. پاسخ در جبر فرمول پایه
است: وقتی $m$ بی‌واسطه در برابری~\ref{eq:desired} بنشیند، $r$ دقیقاً حذف می‌شود.

\begin{equation}
\label{eq:cancel}
d = \left\lceil r \times \frac{m}{T} \right\rceil
  = \left\lceil r \times \frac{\lambda/r}{T} \right\rceil
  = \left\lceil \frac{\lambda}{T} \right\rceil
\end{equation}

یعنی تصمیم سیاست واکنشی اصلاً تابع $r$ نیست، هرچند سیگنالش درون‌زاست. پس درون‌زا بودن سیگنال
به‌خودی‌خود مشکل‌ساز نیست؛ مشکل آنجاست که این حذف‌شدن به هم بخورد.

و پیش‌بینی‌کننده دقیقاً همین کار را می‌کند، چون یک عملگر \emph{حافظه‌دار} است که میان تقسیم
بر $r$ و ضرب در $r$ می‌نشیند: مقدار $\hat{x}_{k+h}$ از تاریخچه‌ی هموارشده‌ی
$m_j = \lambda_j / r_j$ برای $j \le k$ ساخته می‌شود، اما در برابری~\ref{eq:desired} در
$r_k$ \emph{جاری} ضرب می‌شود، و آن $r_j$های گذشته با $r_k$ حذف نمی‌شوند.

اثر را زیر بار کل \emph{ثابت} می‌توان دید. با $\lambda$ ثابت و برابری~\ref{eq:ewmatrend}:

\begin{equation}
\label{eq:gain}
\ell_k - \ell_{k-1} = \alpha\left( \frac{\lambda}{r_k} - \ell_{k-1} \right)
\end{equation}

تنها محرک شیب برآوردشده، تغییر خود $r$ است؛ بار اصلاً تکان نخورده. این شیب سپس در افق $h$
ضرب و به پیش‌بینی افزوده می‌شود، پس $h$ همان دستگیره‌ی بهره‌ی حلقه است و انتظار می‌رود افق
بلندتر ناپایداری را تندتر کند. جدول~\ref{tab:design-oscillation-unit} همین را زیر بار ثابت
نشان می‌دهد.

راه حل، پیش‌بینی بار کل است. \code{PredictiveTotalLoad} نرخ کل ورود درخواست‌ها را از حاصل‌ضرب
$m \times r$ بازسازی می‌کند، همان را پیش‌بینی می‌کند، و ناوگان را چنین اندازه می‌گیرد:

\begin{equation}
\label{eq:fleetfor}
d = \left\lceil \frac{\hat{\lambda}_{k+h}}{T} \right\rceil
\end{equation}

نرخ کل ورود درخواست‌ها یک ورودی \emph{برون‌زا} است، چون مقیاس‌گذار نمی‌تواند تعداد درخواست‌های
کاربران را تغییر دهد. پس بهره‌ی این مسیر ساختاراً صفر است: در
برابری~\ref{eq:fleetfor} هیچ وابستگی‌ای به $r$ دیده نمی‌شود. اینجا ضرب و تقسیم در
\emph{یک چرخه} انجام می‌شوند، پس برخلاف حالت پیش‌بینی بار هر نمونه، هیچ $r$ قدیمی وارد
حاصل نمی‌شود. تنها باقی‌مانده‌ی جزئی از آنجاست که مخرج جمع‌آورنده \code{count(up)} است و
ضریب بازسازی \code{status.readyReplicas}؛ این دو منبع متفاوت‌اند و در گذراها می‌توانند
اندکی از هم فاصله بگیرند. این نکته، تفاوت میان دو سیاست را از یک مشاهده‌ی تجربی به یک ویژگی
ساختاری ارتقا می‌دهد.

اثر را در حلقه‌ی بسته و مستقل از هر زیرساختی اندازه گرفتیم، در آزمون
\code{TestPerReplica\-Forecasting\-Oscillates\-Under\-ConstantLoad}. بار کل ثابت ۵۰۰ درخواست بر ثانیه
بود، هدف ۱۰۰، پیش‌بینی‌کننده \lr{EWMA} با افق ۳ و $\alpha$ برابر ۰/۵. چهل گام اجرا شد که ده
گام نخست آن به عنوان شیب اولیه کنار گذاشته شد. سازوکار پایدارسازی خاموش بود.

\begin{table}[H]
\centering
\caption{نوسان تحت بار کل \emph{ثابت}، در آزمون واحد حلقه‌بسته}
\label{tab:design-oscillation-unit}
\small
\begin{tabular}{|l|c|l|}
\hline
\textbf{سیاست} & \textbf{تغییر تعداد نمونه در ۳۰ گام پایدار} & \textbf{ناوگان} \\
\hline
\code{PredictiveTotalLoad} & \textbf{0} & روی ۵ می‌نشیند و تکان نمی‌خورد \\
\code{PredictivePerReplica} & \textbf{29} & میان ۱ و ۸ نوسان می‌کند \\
\hline
\end{tabular}
\end{table}

خاموش‌بودن سازوکار پایدارسازی در این آزمون عمدی است. آن سازوکار یک میراگر است، و میراکردن
نشانه همان ویژگی‌ای را پنهان می‌کند که موضوع آزمون است. ناپایداری در دینامیک خود سیاست
زندگی می‌کند و نه در خروجی میراشده‌ی آن.

همین پدیده در شبیه‌ساز هم دیده می‌شود، این بار با پیکربندی کامل و بار کاری \arm{spike}:

\begin{table}[H]
\centering
\caption{همان پدیده در شبیه‌ساز، با پیکربندی کامل}
\label{tab:design-oscillation-sim}
\begin{tabular}{|l|c|l|l|}
\hline
\textbf{سیاست} & \textbf{تغییر تعداد نمونه} & \textbf{هزینه} & \textbf{نقض \lr{SLA}} \\
\hline
\code{Threshold} & 4 & مبنا & مبنا \\
\code{PredictivePerReplica} & \textbf{17} & ‎+۲۳ درصد & یکسان \\
\code{PredictiveTotalLoad} & \textbf{3} & برابر \code{Threshold} & بهتر روی \arm{diurnal} \\
\hline
\end{tabular}
\end{table}

بیست‌وسه درصد هزینه‌ی بیشتر، آن هم بدون هیچ سودی در کیفیت خدمت. تیزترین بیان شکست شکل «هر
نمونه» همین است: روی هر دو محور اکیداً بدتر است.

باید تأکید کرد که ارقام جدول‌های~\ref{tab:design-oscillation-unit}
و~\ref{tab:design-oscillation-sim} \emph{قابل تعویض با یکدیگر نیستند}. ارقام آزمون واحد
می‌گویند خود سیاست به‌تنهایی چه می‌کند؛ ارقام شبیه‌ساز می‌گویند پیکربندی کامل چه می‌کند. هر
دو در ارزیابی جای دارند، به شرط آنکه برچسب‌گذاری شوند.

\code{PredictivePerReplica} عمداً در کد باقی مانده. تقابل این دو خودش یک نتیجه‌ی ارزیابی است
و نه کد مرده. برای آنکه به‌اشتباه به‌عنوان سیاست عملیاتی انتخاب نشود، هنگام گزینش آن یک
هشدار در گزارش‌ها ثبت می‌شود.

\section{سازوکار پایدارسازی}
\label{sec:design-stabilizer}

این سازوکار پاسخ پیاده‌سازی به نیازمندی غیرکارکردی «پرهیز از نوسان» است و سه بخش دارد.

\textbf{نخست، پنجره‌ی پایدارسازی نامتقارن.} مقیاس‌گیری به بالا بی‌درنگ اعمال می‌شود، چون کندی
در افزودن ظرفیت مستقیماً به تأخیر کاربر تبدیل می‌شود. مقیاس‌گیری به پایین اما بر پایه‌ی
\emph{بیشینه‌ی} پیشنهادهای دیده‌شده در طول پنجره‌ی $W$ تصمیم گرفته می‌شود:

\begin{equation}
\label{eq:stabilize}
d_{\mathrm{applied}} =
\begin{cases}
d_k & d_k > r \quad\text{(مقیاس به بالا، بی‌درنگ)}\\[2mm]
\max\{\, d_j \;:\; t_k - W \le t_j \le t_k \,\} & d_k < r \quad\text{(مقیاس به پایین)}
\end{cases}
\end{equation}

مقدار پیش‌فرض $W$ برابر ۹۰ ثانیه است. به این ترتیب یک افت کوتاه بار نمی‌تواند ناوگان را کوچک
کند. خود \lr{Kubernetes} نیز در \lr{scaleDown.stabilizationWindowSeconds} همین رفتار را دارد.

\textbf{دوم، شیب‌گیری}\LTRfootnote{Cooldown}. حداقل فاصله‌ی زمانی میان دو کنش مقیاس‌گذاری.
این مؤلفه در نمونه‌ی اولیه نبود و در پیشنهاد پروژه نام برده شده بود. نکته‌ی مهم پیاده‌سازی آن
در بخش~\ref{sec:design-defects} می‌آید: شیب‌گیری باید \emph{به تفکیک جهت} باشد.

\textbf{سوم، سقف گام.} بیشینه‌ی تعداد نمونه‌ای که یک تصمیم منفرد می‌تواند اضافه یا کم کند.

در همه‌ی آزمایش‌های این نوشتار، شیب‌گیری و سقف گام غیرفعال بوده‌اند و تنها پنجره‌ی پایدارسازی
فعال بوده است. دلیل این انتخاب آن است که متغیر مورد مطالعه یکی بماند: هر سه مؤلفه میراگرند و
فعال‌بودن هم‌زمانشان سهم هرکدام را غیرقابل تفکیک می‌کند.

باید صریح گفت که این سازوکار یک \emph{میراگر} است و نه یک \emph{اصلاح}. فصل~\ref{ch:results}
نشان می‌دهد که همین میراگر می‌تواند یک قانون کنترل معیوب را چنان بپوشاند که از همه‌ی آزمون‌ها
سربلند بیرون بیاید. این خودش یکی از نتایج پروژه است.

\section{مقیاس‌گذار}
\label{sec:design-scaler}

مقیاس‌گذار تصمیم را از راه زیرمنبع \lr{scale} روی \lr{Deployment} هدف اعمال می‌کند، و تنها
وقتی که عدد نسبت به \lr{spec.replicas} تغییر کرده باشد.

یک تصمیم پیاده‌سازی در اینجا ارزش ثبت دارد. عملیات \code{Get} از زیرمنبع \lr{scale} استفاده
\emph{نمی‌کند} و به‌جای آن یک بار خودِ \lr{Deployment} را می‌خواند. دلیلش آن است که زیرمنبع
\lr{scale} نمی‌تواند به پرسش دوم پاسخ دهد: وضعیت آن تعداد \emph{کل} نمونه‌ها را گزارش می‌کند
و نه تعداد نمونه‌های \emph{آماده} را، حال آنکه فرمول به عدد دوم نیاز دارد. خواندن یک‌باره‌ی
\lr{Deployment} هر دو عدد را می‌دهد و تعداد فراخوانی واسط برنامه‌نویسی را در هر چرخه نصف
می‌کند. عملیات \code{Set} همچنان از \lr{GetScale} و \lr{UpdateScale} استفاده می‌کند تا
چرخه‌ی خواندن-تغییر-نوشتن روی \lr{resourceVersion} درست انجام شود.

پیاده‌سازی دوم این واسط، یعنی \code{InMemory}، همان است که آزمون‌های حلقه‌بسته‌ی
بخش~\ref{sec:design-perreplica} را ممکن می‌کند.

\section{مشاهده‌پذیری}
\label{sec:design-observability}

مقیاس‌گذار سنجه‌های داخلی خودش را روی یک نقطه‌ی \lr{/metrics} منتشر می‌کند. مهم‌ترین آن‌ها در
جدول~\ref{tab:design-metrics} آمده است.

\begin{table}[H]
\centering
\small
\caption{سنجه‌های داخلی منتشرشده توسط مقیاس‌گذار}
\label{tab:design-metrics}
\begin{tabular}{|l|l|}
\hline
\textbf{نام سنجه} & \textbf{معنا} \\
\hline
\code{autoscaler\_current\_replicas} & تعداد نمونه‌ی درخواست‌شده (\lr{spec.replicas}) \\
\code{autoscaler\_ready\_replicas} & تعداد نمونه‌ی آماده؛ پایه‌ی فرمول \\
\code{autoscaler\_desired\_replicas} & تصمیم سیاست، پس از پایدارسازی \\
\code{autoscaler\_raw\_recommendation} & پیشنهاد خام، پیش از پایدارسازی \\
\code{autoscaler\_metric\_value} & مقدار مشاهده‌شده‌ی سیگنال \\
\code{autoscaler\_metric\_target} & نقطه‌ی هدف \\
\code{autoscaler\_predicted\_value} & مقدار پیش‌بینی‌شده (صفر برای سیاست‌های واکنشی) \\
\code{autoscaler\_reconcile\_duration\_seconds} & زمان اجرای هر چرخه \\
\code{autoscaler\_scale\_actions\_total} & شمار کنش‌های مقیاس‌گذاری \\
\code{autoscaler\_collector\_errors\_total} & خطاهای خواندن سنجه \\
\hline
\end{tabular}
\end{table}

این سنجه‌ها سه کاربرد دارند. یکی مصورسازی در \lr{Grafana}. دوم، پایش سلامت خود مقیاس‌گذار.
سوم \mdash{} و مهم‌ترین برای این نوشتار \mdash{} یک \emph{بررسی اعتبار مستقل}. تفاوت میان
\code{autoscaler\_raw\_recommendation} و \code{autoscaler\_desired\_replicas} دقیقاً می‌گوید
میراگر چه مقدار کار انجام داده است؛ و مقایسه‌ی \code{autoscaler\_metric\_value} با بار
مشاهده‌شده‌ی هر نمونه، امضای پس‌خور مثبت را آشکار می‌کند. در
بخش~\ref{sec:results-stabilizer} از هر دو استفاده شده است.

یک تصمیم کوچک اما آگاهانه: نام سنجه‌ها پیشوند \code{autoscaler\_} را نگه داشته‌اند، هرچند نام
پروژه به \lr{setpoint} تغییر کرده است. تغییر نام یک سنجه، هر پنل \lr{Grafana} و هر شکل
فصل ارزیابی را بی‌اعتبار می‌کند. یک آزمون این نام‌ها را تثبیت کرده است.

\section{حلقه‌ی کنترل}
\label{sec:design-loop}

هر \lr{interval\_seconds} یک چرخه‌ی آشتی‌دهی اجرا می‌شود. مقدار پیش‌فرض ۱۵ ثانیه است که با
دوره‌ی همگام‌سازی \lr{HPA} هم‌اندازه است. هر چرخه شش گام دارد:

\begin{enumerate}
  \item خواندن سیگنال مقیاس‌گذاری از \lr{Prometheus} با \lr{PromQL}.
  \item خواندن \lr{spec.replicas} و \lr{status.readyReplicas} از هدف.
  \item گرفتن پیشنهاد خام از سیاست.
  \item پایدارسازی: مقیاس به بالا بی‌درنگ، مقیاس به پایین مشروط به بیشینه‌ی پنجره، سپس سقف
  گام و شیب‌گیری.
  \item اعمال از راه زیرمنبع \lr{scale}، در صورت تغییر عدد.
  \item صدور تصمیم و ورودی‌هایش روی \lr{/metrics}.
\end{enumerate}

شکست در هر گام ثبت می‌شود و تعداد نمونه‌ی جاری نگه داشته می‌شود. این انتخاب \mdash{} «نگه‌داشتن» به
جای «حدس زدن» \mdash{} تنها پاسخ امن به نبود اطلاعات است: قطعی \lr{Prometheus} شاهدی بر بالا رفتن یا
پایین آمدن بار نیست.

\textbf{سربار خود مقیاس‌گذار.} در یک اجرای واقعی، ۱۸ چرخه‌ی آشتی‌دهی مجموعاً ۰/۰۶۸ ثانیه طول
کشید، یعنی حدود ۳/۸ میلی‌ثانیه بر چرخه در برابر یک فاصله‌ی ۱۵ ثانیه‌ای. حلقه‌ی کنترل به هیچ
وجه گلوگاه نیست، و این عدد پاسخ روشنی به پرسش «آیا زمان واکنش اندازه‌گیری‌شده‌ی شما صرفاً
کندی کنترل‌گر خودتان نیست؟» می‌دهد. این همان نیازمندی غیرکارکردی «سربار منابع پایین» است.

یک حالت \lr{dry-run} نیز پیاده‌سازی شده که تصمیم‌ها را ثبت می‌کند اما اعمال نمی‌کند. کاربرد
آن، اجرای مقیاس‌گذار در کنار یک سرویس واقعی برای دیدن اینکه \emph{چه می‌کرد} است، پیش از آنکه
اجازه‌ی نوشتن بگیرد.

\section{پیکربندی}
\label{sec:design-config}

کل رفتار سامانه از یک فایل \lr{YAML} می‌آید که در خوشه به شکل \lr{ConfigMap} نصب می‌شود. طرح
این فایل عمداً با فایل پیکربندی نمونه‌ی اولیه‌ی پایتونی یکی است. به این ترتیب مقایسه‌ی
شبیه‌ساز و خوشه در فصل~\ref{ch:results} با تفاوت تنظیمات مخدوش نمی‌شود.

جدول~\ref{tab:design-config} مقادیری را نشان می‌دهد که در همه‌ی اندازه‌گیری‌ها به کار رفت.
فایل کامل در پیوست الف آمده است.

\begin{table}[H]
\centering
\small
\caption{پیکربندی به‌کاررفته در همه‌ی اجراهای ارزیابی}
\label{tab:design-config}
\begin{tabular}{|l|c|l|}
\hline
\textbf{کلید} & \textbf{مقدار} & \textbf{توضیح} \\
\hline
\code{interval\_seconds} & 15 & هم‌اندازه‌ی دوره‌ی همگام‌سازی \lr{HPA} \\
\code{policy.target} & 100.0 & درخواست بر ثانیه بر هر نمونه‌ی آماده \\
\code{policy.tolerance} & 0.10 & ناحیه‌ی مرده‌ی ۱۰ درصدی، مانند \lr{HPA} \\
\code{policy.min\_replicas} & 1 & \\
\code{policy.max\_replicas} & 20 & بنگرید به توضیح زیر \\
\code{forecaster.type} & \code{ewma} & \code{EWMATrend} \\
\code{forecaster.horizon} & 3 & ‏۳ × ۱۵ ثانیه = ۴۵ ثانیه پیش‌دستی \\
\code{forecaster.alpha} & 0.5 & \\
\code{stabilization\_window\_seconds} & 90 & صفر در بازوهای حذفی \\
\code{cooldown\_seconds} & 0 & غیرفعال \\
\code{max\_step} & 0 & غیرفعال \\
\hline
\end{tabular}
\end{table}

انتخاب عدد ۲۰ برای بیشینه‌ی نمونه خودش یک تصمیم ارزیابی است و نه یک عدد دلبخواه. همه‌ی چهار
الگوی بار در اوج به ۱۲۰۰ تا ۱۳۰۰ درخواست بر ثانیه می‌رسند، که در برابر هدف ۱۰۰ یعنی ۱۲ تا ۱۳
نمونه. اگر سقف ۱۰ می‌بود، \emph{همه‌ی} بازوهای خودکار در اوج به سقف می‌خوردند و هر کدام با ۱۰
نمونه و ۱۳۰ درخواست بر ثانیه بر نمونه کار می‌کردند \mdash{} یعنی همگی یکسان و بالای خط کیفیت خدمت.
آنگاه سنجه‌ی نقض کیفیت خدمت دقیقاً در همان نقطه‌ای که مقایسه بیشترین اهمیت را دارد، بازوها را
از هم تفکیک نمی‌کرد، و سیاست پیش‌بینانه به دلیلی که هیچ ربطی به سیاست ندارد بهتر از واکنشی به
نظر نمی‌رسید. ظرفیت گره محدودیت نبود: ۲۰ نمونه ۴ هسته از ۱۴ هسته‌ی گره را درخواست می‌کنند.

\section{نقص‌هایی که آزمون‌ها آشکار کردند}
\label{sec:design-defects}

سه نقص در مسیر کنترل با نوشتن آزمون پیدا شد و نه با خواندن کد. هر سه وقتی آشکار شدند که
رفتار مورد انتظار را ادعا کردیم و ادعا شکست خورد. هیچ‌کدام با مطالعه‌ی کد دیده نمی‌شد، و
نکته‌ی مشترک هر سه آن است که تنها در بدترین وضعیت خوشه فعال می‌شوند.

\textbf{نخست، \code{Holt.Update} روی مسیر \lr{NaN} پیش‌بینی منفی برمی‌گرداند.} بازگشت زودهنگام
برای \lr{NaN} و \lr{Inf}، کف \lr{math.Max(0, ...)} را رد می‌کرد که مسیر عادی اعمال می‌کند. با
تغذیه‌ی ۵۰۰، ۴۰۰، ۱۰۰، ۱۰ و سپس \lr{NaN} با افق ۳ و $\alpha = \beta = 0.5$، مقدار منفی ۴۵۰
اندازه‌گیری شد. \code{FleetFor} این عدد را بی‌سروصدا به \lr{min\_replicas} کف می‌زد. یعنی یک
مقیاس‌گیری به پایین که علتش یک \lr{scrape} ازدست‌رفته بود و نه نبود بار. در گزارش‌ها هم
کاملاً عادی به نظر می‌رسید.

\textbf{دوم، شیب‌گیری سراسری بود و نه به تفکیک جهت.} این رفتار با توضیح خود تابع در تناقض
بود، که گفته بود مقیاس به بالا بی‌درنگ اعمال می‌شود. با طرح پروژه هم در تناقض بود. پیامدش
این بود که پس از یک مقیاس به پایین، بازگشت بار نمی‌توانست مقیاس به بالای جبرانی را در تمام
مدت شیب‌گیری فعال کند. یعنی دقیقاً همان حالتی که انتظار کشیدن در آن به قیمت نقض کیفیت خدمت
تمام می‌شود. اکنون بر پایه‌ی جهت کلید خورده است.

\textbf{سوم، تبدیل \lr{int32(math.Ceil(...))} به رفتار وابسته به پیاده‌سازی تکیه داشت.}
مشخصه‌ی زبان \lr{Go} می‌گوید تبدیل اعشاری به صحیح، وقتی نتیجه در نوع مقصد نگنجد، موفق می‌شود
اما مقدار نتیجه وابسته به پیاده‌سازی است. پس یک پیشنهاد بسیار بزرگ ممکن است روی یک معماری به
\lr{MaxInt32} اشباع شود و روی معماری دیگر به عددی \emph{منفی} بپیچد. آنگاه \code{clamp} عدد
منفی را «کمتر از کمینه» می‌خواند و ناوگان را زیر بار شدید \emph{کوچک} می‌کند. بدترین پاسخ
ممکن، که تصادفی به دست آمده باشد. آزمون‌ها روی این میزبان \lr{arm64} در هر دو حالت سبز
می‌مانند و دقیقاً به همین دلیل یک نگهبان صریح لازم بود، چون کانتینر برای معماری‌ای ساخته
می‌شود که با معماری توسعه یکی نیست. تابع \code{ceilToReplicas} اکنون عمداً اشباع می‌کند.

\section{برنامه‌ی نمونه و بستر آزمایش}
\label{sec:design-sampleapp}

برای آنکه ارزیابی معنا داشته باشد، سرویس هدف باید واقعاً محدود به پردازنده باشد. سرویسی که
به‌جای محاسبه فقط بخوابد، زیر بار فزاینده مصرف پردازنده‌ی ثابتی نشان می‌دهد. آنگاه بازوی
\arm{hpa-cpu} هرگز واکنش نشان نمی‌دهد و خط مبنا به یک مترسک تبدیل می‌شود.

پس برنامه‌ی نمونه برای هر درخواست مقدار مشخصی پردازنده می‌سوزاند. این مقدار چنان کالیبره شده
که نمونه‌ای در نرخ هدف خود دقیقاً ۱۰۰ درصد درخواست پردازنده‌اش را مصرف کند. جزئیات این
کالیبراسیون و تأیید آن در خوشه در بخش~\ref{sec:method-calibration} آمده است.

برنامه‌ی نمونه سه سنجه منتشر می‌کند: شمارنده‌ی \code{http\_requests\_total}، بافه‌ی
\code{http\_request\_\-duration\_\-seconds}، و سنجه‌ی لحظه‌ای \code{http\_requests\_in\_flight}.
دو پارامتر محیطی نیز دارد که رفتار آن را تنظیم می‌کنند: \code{WORK\_CPU\_MS} برای مقدار
پردازش هر درخواست و \code{WORK\_MAX\_CONCURRENT} برای سقف هم‌زمانی.

بستر آزمایش شامل \lr{Prometheus} و \lr{Grafana} از چارت \lr{kube-prometheus-stack} است،
به‌علاوه‌ی یک \lr{ServiceMonitor} برای سرویس هدف و مولد بار \lr{k6}. در این بستر چند دام
خاموش وجود داشت که کشفشان زمان برد و ثبتشان ارزش دارد.

\textbf{\lr{targetLabels: [app]} روی \lr{ServiceMonitor} الزامی است.} عبارت \lr{PromQL} روی
برچسب \lr{app="sample"} انتخاب می‌کند، اما این برچسب روی \lr{Service} است و نه در سنجه‌هایی
که خود برنامه منتشر می‌کند. بدون ارتقای این برچسب، پرس‌وجو هیچ‌چیز را تطبیق نمی‌دهد، برداری
تهی برمی‌گرداند، و ناوگان روی \lr{min\_replicas} می‌نشیند. همه‌ی این‌ها هم کاملاً خاموش رخ
می‌دهد.

\textbf{\lr{serviceMonitorSelectorNilUsesHelmValues} به‌صورت پیش‌فرض \lr{true} است.} در این
حالت \lr{Prometheus} فقط \lr{ServiceMonitor}هایی را کشف می‌کند که برچسب انتشار \lr{Helm} را
دارند. نمونه‌های دست‌نویس بی‌هیچ توضیحی در هیچ گزارشی نادیده گرفته می‌شوند.

\textbf{\code{http\_requests\_total} تا پیش از نخستین درخواست وجود ندارد.} نتیجه‌ی تهی
پرس‌وجو پیش از هر ترافیکی، رفتار درست است و نه خط لوله‌ی خراب. همین نکته یکی از دلایل وجود
مرحله‌ی گرم‌کردن در چارچوب آزمایش است (بخش~\ref{sec:method-harness}).

\textbf{\lr{kubectl port-forward} نمی‌تواند به یک ناوگان بار بدهد.} این فرمان \emph{یک} نقطه‌ی
انتهایی را حل می‌کند و همان را نگه می‌دارد. هر آزمایشی از این راه، به ظرفیت یک نمونه محدود
می‌ماند در حالی که ناوگان زیر آن بزرگ می‌شود و بی‌استفاده می‌ماند \mdash{} و همه‌ی بازوها یکسان به
نظر می‌رسند چون هیچ‌کدام واقعاً بارگذاری نشده‌اند. به همین دلیل \lr{Service} از نوع
\lr{LoadBalancer} تعریف شده است.

\section{جمع‌بندی فصل}
\label{sec:design-summary}

سامانه در امتداد سه دغدغه‌ی مستقل شکسته شده و هر مرحله واسطی با دست‌کم دو پیاده‌سازی دارد.
این ساختار، آزمون حلقه‌بسته‌ی سیاست را بدون خوشه ممکن می‌کند و یافته‌ی مرکزی پروژه از همین
راه به دست آمد.

دو انحراف از نمونه‌ی اولیه هر دو با اندازه‌گیری توجیه شده‌اند. پایه‌ی فرمول باید نمونه‌های
آماده باشد. پیش‌بینی باید روی بار کل انجام شود و نه بار هر نمونه، و دلیل آن ساختاری است:
اندازه‌ی ناوگان در برابری~\ref{eq:fleetfor} به تعداد نمونه‌ی جاری وابسته نیست.

سازوکار پایدارسازی یک میراگر است و نه اصلاح. فصل~\ref{ch:results} نشان می‌دهد چه چیزی را
می‌پوشاند.

سه نقص مسیر کنترل تنها با نوشتن آزمون به‌مثابه‌ی مشخصه آشکار شدند. هر سه در شرایطی فعال
می‌شدند که خوشه در بدترین وضعیت خود است: یک \lr{scrape} ازدست‌رفته، بازگشت بار پس از یک
مقیاس به پایین، و بار شدید روی معماری دیگر.
